% chapters/13/section-04.tex
\section{Trigger Volume dan Interaksi}
\label{sec:trigger}

\subsection{Kode Lengkap}

\kodecsharp{unity/ch13_trigger_pintu.cs}{Trigger volume yang menandai pemain telah mencapai suatu tempat}{lst:ch13-trigger}

\subsection{Fungsi yang Dipanggil Ketika Terjadi Tabrakan}

Skrip itu tidak memeriksa tabrakan sendiri. Ia hanya
menyediakan dua fungsi, yaitu \code{OnTriggerEnter2D} dan
\code{OnTriggerExit2D}, yang akan dipanggil mesin ketika peristiwanya terjadi.

Cara itu berbeda dari yang Anda lakukan pada Bab~VII, tempat Anda memeriksa
tabrakan sendiri melalui \code{spritecollide}. Pada Unity, pemeriksaannya
diserahkan kepada mesin, dan kode Anda hanya menanggapi hasilnya. Inilah
padanan dari pembedaan ``mendeteksi'' dan ``menyelesaikan'' yang telah
dijelaskan pada Bagian~7.2: mesin mendeteksi, sedangkan Anda menentukan apa
yang terjadi.

\subsection{Syarat agar Trigger Bekerja}

Tiga syarat berikut harus dipenuhi agar fungsi trigger dipanggil. Ketiganya
sering terlupakan, dan gejalanya sama, yaitu fungsi yang tidak pernah
dijalankan.

\begin{enumerate}
  \item \textbf{\istilah{Collider} benda itu harus ditandai sebagai
        \istilah{Is Trigger} pada \istilah{Inspector}.}
  \item \textbf{Sekurang-kurangnya satu dari kedua benda harus memiliki
        \istilah{Rigidbody}.} Pemain biasanya sudah memilikinya, sehingga
        syarat ini jarang menjadi penyebab.
  \item \textbf{Kedua benda harus berada pada \istilah{layer} yang saling
        diperiksa.} Pengaturan \istilah{layer} menentukan pasangan benda mana
        yang tabrakannya diperiksa mesin, dan sering dipakai untuk
        menghemat perhitungan.
\end{enumerate}

\begin{catatabox}
Perhatikan pemakaian peubah \code{alreadyTriggered} pada skrip di atas. Peubah
itu mencegah pesan tercetak berulang kali selama pemain berada di dalam daerah
trigger. Tanpa peubah itu, fungsinya tetap dipanggil sekali saja, tetapi isinya
dapat dijalankan berkali-kali apabila terdapat beberapa \istilah{Collider} pada
benda yang sama. Kebiasaan memeriksa keadaan sebelum bertindak berlaku juga di
sini.
\end{catatabox}

\subsection{Tag dan Layer}

Dua pengaturan yang berkaitan dengan pengelompokan benda perlu Anda kenal.

\begin{description}
  \item[\istilah{Tag}] Penanda pada sebuah benda, misalnya \code{Player} atau
        \code{Enemy}. Dipakai oleh kode Anda untuk mengenali benda, misalnya
        pada baris \code{other.CompareTag(playerTag)}.
  \item[\istilah{Layer}] Pengelompokan benda untuk keperluan pemeriksaan
        tabrakan oleh mesin. Dipakai untuk menentukan pasangan benda mana yang
        diperiksa, misalnya agar peluru tidak memeriksa tabrakan dengan sesama
        peluru.
\end{description}

Perbedaan keduanya mudah dikacaukan. Cara mengingatnya: \istilah{tag} dipakai
oleh \textbf{kode Anda}, sedangkan \istilah{layer} dipakai oleh
\textbf{mesin}.
